\par\Needspace{4\baselineskip}
\begin{center}
\hypertarget{native-035-article017}{}\pdfbookmark[2]{17.}{article-017}17.
\end{center}
Zur Bestimmung des \emph{Maximums der Mittelpunktsgleichung} seien noch einige Bemerkungen hinzugefügt. Zunächst ist offensichtlich, dass der Unterschied zwischen der exzentrischen und der mittleren Anomalie bei $E=90^0$ am größten ist und dort $=e$ ist (ausgedrückt in Grad und deren Unterteilungen). Der Radiusvektor ist an diesem Punkt $=a$, woraus sich $v=90^0+\varphi$ und damit für die gesamte Mittelpunktsgleichung $=\varphi+e$ ergibt. Sie erreicht hier noch nicht ihr Maximum, da der Unterschied zwischen $v$ und $E$ noch über $\varphi$ hinauswachsen kann. \emph{Dieser} Unterschied ist am größten, wenn $\mathrm{d}(v-E)=0$ oder $\mathrm{d}v=\mathrm{d}E$ gilt. Dabei ist die Exzentrizität offenbar als konstant anzusehen. Unter dieser Voraussetzung gilt allgemein $\frac{\mathrm{d}v}{\sin v}=\frac{\mathrm{d}E}{\sin E}$. Also muss an dem Punkt, an dem der Unterschied zwischen $v$ und $E$ am größten ist, $\sin v=\sin E$ gelten. Aus den Gleichungen VIII und III folgt daher
\[
r=a\cos\varphi
\]
\[
e\cos E=1-\cos\varphi
\]
oder
\[
\cos E=+\operatorname{tang}\frac{1}{2}\varphi.
\]
Entsprechend findet man $\cos v=-\operatorname{tang}\frac{1}{2}\varphi$, also gilt\footnote{Die Maxima zwischen Aphel und Perihel können außer Betracht bleiben, da sie sich von denen zwischen Perihel und Aphel offenbar nur im Vorzeichen unterscheiden.}
\[
v=90^0+\operatorname{arc.}\sin\operatorname{tang}\frac{1}{2}\varphi
\]
\[
E=90^0-\operatorname{arc.}\sin\operatorname{tang}\frac{1}{2}\varphi;
\]
Daraus folgt weiter $\sin E=\BareRoot{(1-\operatorname{tang}\frac{1}{2}\varphi^2)}=\frac{\BareRoot{\cos\varphi}}{\cos\frac{1}{2}\varphi}$, sodass die gesamte Mittelpunktsgleichung an diesem Punkt
\[
=2\operatorname{arc.}\sin\operatorname{tang}\frac{1}{2}\varphi+2\sin\frac{1}{2}\varphi\cdot\BareRoot{\cos\varphi}
\]
beträgt, wobei der zweite Term in Grad und deren Unterteilungen auszudrücken ist.

An dem Punkt schließlich, an dem die gesamte Mittelpunktsgleichung selbst ihr Maximum erreicht, muss $\mathrm{d}v=\mathrm{d}M$ und damit nach Art. 15 $r=a\BareRoot{\cos\varphi}$ gelten. Daraus folgt
\SourceLeaf{036}
\[
\cos v=-\frac{1-\cos\varphi^{\frac{3}{2}}}{e}
\]
\[
\cos E=\frac{1-\BareRoot{\cos\varphi}}{e}=\frac{1-\cos\varphi}{e(1+\BareRoot{\cos\varphi})}=\frac{\operatorname{tang}\frac{1}{2}\varphi}{1+\BareRoot{\cos\varphi}},
\]
woraus sich $E$ mit höchster Genauigkeit bestimmen lässt. Ist $E$ gefunden, so ergibt sich nach den Gleichungen X und XII für die Mittelpunktsgleichung $=2\operatorname{arc.}\sin\frac{\sin\frac{1}{2}\varphi\sin E}{\BareRoot[4]{\cos\varphi}}+e\sin E$\HistoricalNoteEight{Handschriftliche Bemerkung:] $\sin E=\BareRoot{\frac{\cos\varphi+2\BareRoot{\cos\varphi}+\cos\varphi^{\frac{3}{2}}}{(1+\BareRoot{\cos\varphi})(1+\cos\varphi)}},$ Mittelpunktsgleichung $=2\operatorname{arc.}\sin\BareRoot{\left(\frac{1}{1+\cos\varphi}-\frac{\BareRoot[4]{(1-e e)}}{2}\right)}+e\sin E=\operatorname{arc.}\cos\frac{-1+\BareRoot{\cos\varphi}+\cos\varphi+\cos\varphi^{\frac{3}{2}}}{1+\cos\varphi}+e\sin E=\operatorname{arc.}\sin(\operatorname{tg}\frac{1}{2}\varphi^2\operatorname{tg}E)+e\sin E.$}. Auf die Darstellung des Maximums der Mittelpunktsgleichung als Reihe in Potenzen der Exzentrizität, die mehrere Autoren angegeben haben, gehen wir hier nicht näher ein. Als Beispiel fügen wir eine Übersicht über die drei hier betrachteten Maxima für Juno bei, wobei wir nach den neuesten Elementen die Exzentrizität $=0{,}255\,4996$ ansetzen.
\begin{center}
\footnotesize
\setlength{\tabcolsep}{3pt}
\renewcommand{\arraystretch}{1.3}
\begin{tabular}{|c|c|c|c|c|}
\hline
Maximum & $E\quad\text{·}$ & $E-M$ & $v-E$ & $v-M$\\
\hline
$E-M$ & $90^0\ 0^{\prime}\ 0^{\prime\prime}$ & $14^0\,38^{\prime}\,20^{\prime\prime}{,}57$ & $14^0\,48^{\prime}\,11^{\prime\prime}{,}48$ & $29^0\,26^{\prime}\,32^{\prime\prime}{,}05$\\
$v-E$ & $82\ \ 32\ \ 9$ & $14\ \ 30\ \ 54{,}01$ & $14\ \ 55\ \ 41{,}79$ & $29\ \ 26\ \ 35{,}80$\\
$v-M$ & $86\ \ 14\ \ 40$ & $14\ \ 36\ \ 27{,}39$ & $14\ \ 53\ \ 49{,}57$ & $29\ \ 30\ \ 16{,}96$\\
\hline
\end{tabular}
\end{center}
